% ==========================================================
%  EXAM QUESTIONS TEMPLATE — AI / ML / MATHEMATICS
%  Design Philosophy:
%    Shares the visual language of summary.tex (same accent,
%    same box style, same typography) but uses article class
%    — no chapters, just sections and question boxes.
%
%  HOW TO USE:
%    1. Copy this file into the course folder: exams.tex
%    2. Set \coursetitle below.
%    3. Add \section{} per topic, questions inside
%       \questiontitle{…} … \closequestion blocks.
% ==========================================================

\documentclass[11pt,oneside]{article}

\usepackage[
  a4paper,
  top=2.4cm,
  bottom=2.1cm,
  left=2.3cm,
  right=2.3cm,
  headheight=15pt
]{geometry}

\usepackage[T1]{fontenc}
\usepackage[utf8]{inputenc}
\usepackage{lmodern}
\usepackage{microtype}
\usepackage{inconsolata}
\usepackage{amsmath,amssymb,mathtools,bm}
\usepackage{enumitem}
\usepackage{booktabs,array,multicol,longtable}
\usepackage{xcolor}
\usepackage[most]{tcolorbox}
\usepackage{titlesec}
\usepackage{fancyhdr}
\usepackage[hidelinks]{hyperref}

\setlength{\parindent}{0pt}
\setlength{\parskip}{0.55em}
\setlist[itemize]{topsep=3pt,itemsep=2pt,leftmargin=2.2em}
\setlist[enumerate]{topsep=3pt,itemsep=2pt,leftmargin=2.2em}

% ── Colors ───────────────────────────────────────────────
% Shared accent with summary.tex — do not change.
\definecolor{accent}{HTML}{1B3A5C}
\definecolor{q@frame}{HTML}{1A5276}
\definecolor{q@bg}{HTML}{F6FAFD}
\definecolor{key@frame}{HTML}{1E8449}
\definecolor{key@bg}{HTML}{F4FBF6}
\definecolor{note@frame}{HTML}{9A6A00}
\definecolor{note@bg}{HTML}{FEF9E7}

% ── Box global style ─────────────────────────────────────
\tcbset{
  enhanced,
  breakable,
  arc=2.5mm,
  boxrule=0.5pt,
  left=4.5mm,
  right=4mm,
  top=2.5mm,
  bottom=2.5mm,
  fonttitle=\bfseries\small,
  sharp corners=west,
  parbox=false
}

% Question box (blue) — one per exam question.
\newtcolorbox{qbox}[1][]{
  colback=q@bg,
  colframe=q@frame,
  borderline west={3.5pt}{0pt}{q@frame},
  title={#1}
}

% Key insight box (green) — answer hints, key results.
\newtcolorbox{keybox}[1][]{
  colback=key@bg,
  colframe=key@frame,
  borderline west={3.5pt}{0pt}{key@frame},
  title={#1}
}

% Organizational note box (amber) — structural remarks only.
\newtcolorbox{notebox}[1][]{
  colback=note@bg,
  colframe=note@frame,
  borderline west={3.5pt}{0pt}{note@frame},
  title={#1}
}

% ── Section headings ─────────────────────────────────────
\titleformat{\section}
  {\normalfont\Large\bfseries\color{accent}}
  {}{0em}{}
  [{\vspace{1pt}\color{accent!45}\titlerule[0.45pt]}]
\titlespacing*{\section}{0pt}{3.2ex plus 1ex minus .2ex}{1.8ex plus .2ex}

\titleformat{\subsection}
  {\normalfont\normalsize\bfseries\color{accent!80!black}}
  {}{0em}{}
\titlespacing*{\subsection}{0pt}{2.2ex plus 1ex minus .2ex}{1.0ex plus .2ex}

% ── Header & footer ──────────────────────────────────────
\pagestyle{fancy}
\fancyhf{}
\fancyhead[L]{\small\color{gray!70}Old Exam Questions — \coursetitle}
\fancyhead[R]{\small\color{gray!70}\thepage}
\renewcommand{\headrulewidth}{0.4pt}

% ── Document metadata ─────────────────────────────────────
\newcommand{\coursetitle}{[Course Title Here]}

% ── Question counter & macros ────────────────────────────
\newcounter{question}
\newcommand{\questiontitle}[1]{%
  \stepcounter{question}\begin{qbox}[Question \thequestion: #1]}
\newcommand{\closequestion}{\end{qbox}}

% Mark questions that appeared on multiple exams.
\newcommand{\repeatnote}{\textit{\color{gray!70}Asked multiple times.}}

% Multiple-choice checkbox bullet.
\newcommand{\choice}{\item[$\square$] }

% Fill-in answer line. Usage: \answerline{3cm}
\newcommand{\answerline}[1]{\makebox[#1]{\hrulefill}}

% Blank working area — dashed-border box for student calculations.
% \workingspace{4cm}   short calculation / derivation
% \workingspace{7cm}   open-ended / long answer
\newcommand{\workingspace}[1]{%
  \par\vspace{2pt}%
  \begin{tcolorbox}[
    enhanced,
    breakable=false,
    colback=white,
    colframe=gray!40,
    boxrule=0.45pt,
    arc=0pt,
    sharp corners,
    height=#1,
    left=3mm, right=3mm, top=3mm, bottom=3mm,
    frame style={dashed}
  ]%
  \end{tcolorbox}%
  \vspace{2pt}%
}

% ── Math macros (minimal subset for exam questions) ──────
% Add from template.tex as needed.
\newcommand{\R}{\mathbb{R}}
\newcommand{\E}{\mathbb{E}}
\newcommand{\Prob}{\mathbb{P}}
\newcommand{\KL}[2]{D_{\mathrm{KL}}\!\left(#1\,\|\,#2\right)}
\newcommand{\dd}{\,\mathrm{d}}


% ==========================================================
%  BEGIN DOCUMENT
% ==========================================================
\begin{document}

% ── Title page ───────────────────────────────────────────
\thispagestyle{empty}
\vspace*{1.4cm}
\begin{center}
  {\color{accent}\rule{\linewidth}{1.8pt}}\par
  \vspace{0.65cm}
  {\fontsize{26}{32}\selectfont\bfseries\color{accent}Old Exam Questions\\[2pt] \coursetitle\par}
  \vspace{0.55cm}
  {\color{accent!35}\rule{\linewidth}{0.5pt}}\par
  \vspace{0.55cm}
  {\large Stefan Eder\par}
  \vspace{0.15cm}
  {\small\color{gray!70}\today\par}
\end{center}

\vspace{0.9cm}
\begin{notebox}[How this file is organized]
This file is separate from the course summary and is meant as a clean practice set.
Exact duplicates were removed. When essentially the same question appeared in multiple old exams,
that is marked with \repeatnote\ inside the question box.
The wording was lightly normalized where the scans were messy, but the mathematical content
and answer choices were kept intact.
\end{notebox}

\tableofcontents
\newpage


% ==========================================================
%  CONTENT — one \section{} per topic area
%  Use \questiontitle{Short title} … \closequestion for each question.
% ==========================================================

\section{[First Topic Area]}

% --- Multiple-choice ---
\questiontitle{[MCQ title]}
[Question text here.]
\begin{itemize}[leftmargin=1.8em]
  \choice Option A
  \choice Option B
  \choice Option C
  \choice Option D
\end{itemize}
\closequestion

% --- Single number / short expression ---
\questiontitle{[Fill-in title]}
[Question text.] \answerline{3cm}
\closequestion

% --- Calculation / short derivation ---
\questiontitle{[Derivation title]}
[Question text.]
\workingspace{4cm}
\closequestion

% --- Open / long answer ---
\questiontitle{[Open question title]}
[Question text.]
\workingspace{7cm}
\closequestion


% ==========================================================
%  ANSWER KEY — always the last section; never add answers
%  inside \questiontitle / \closequestion blocks above.
% ==========================================================
\newpage
\section*{Answer Key}

\begin{enumerate}[label=\textbf{Q\arabic*.}, leftmargin=3em, itemsep=4pt]
  \item \textbf{B}
  \item $x = 0.5$
  \item [key steps: \dots]
  \item [key points: \dots]
\end{enumerate}


\end{document}
